\chapter{Mercado Objetivo}

\section{Identificación del mercado objetivo}

El mercado objetivo de MuseIQ está conformado por instituciones culturales que desarrollan recorridos presenciales y que requieren fortalecer sus mecanismos de mediación, orientación y accesibilidad mediante herramientas digitales. En términos operativos, la propuesta se dirige a museos, museos de sitio, centros culturales, espacios de interpretación patrimonial y otras entidades que administran colecciones, exposiciones o circuitos de visita donde el usuario necesita información contextual durante el desplazamiento.

La identificación de este mercado se sustenta en dos condiciones principales. La primera es la existencia de una base institucional real para la adopción. El Ministerio de Cultura del Perú administra una red de 56 museos a nivel nacional, lo que evidencia un ecosistema museístico suficientemente amplio para plantear una solución replicable y no restringida a un único caso de uso~\cite{museoscultura}. La segunda condición es la disponibilidad creciente de dispositivos móviles entre la población peruana. OSIPTEL reportó que, en 2024, el 94.8\% de los hogares peruanos contaba con al menos un smartphone, lo que reduce la necesidad de dotar a cada visitante de hardware dedicado y favorece arquitecturas basadas en el dispositivo personal del usuario~\cite{osiptel2024smartphone}.

Dentro de este universo institucional, los primeros candidatos de adopción son aquellos espacios con afluencia sostenida, recorridos autoguiados o semiguiados, y necesidad de enriquecer la experiencia interpretativa sin incrementar de forma proporcional la carga de personal. En este grupo destacan los museos estatales con alta visibilidad pública, los museos de sitio asociados a patrimonio arqueológico, los museos universitarios y privados con interés en innovación, así como centros culturales que desarrollan exposiciones temporales o permanentes. Todos ellos comparten un problema funcional: la información relevante para el visitante suele depender de cartelas, paneles, audioguías lineales o mediación humana limitada en cobertura temporal.

En ese contexto, una guía inteligente como MuseIQ resulta pertinente porque combina localización indoor mediante BLE, interacción por voz y recuperación contextual de contenidos para ofrecer acompañamiento durante el recorrido. La propuesta no sustituye la mediación humana especializada, pero sí amplía su alcance, mejora la accesibilidad y permite entregar información contextualizada en el momento y lugar de la visita. Esta lógica de apoyo resulta consistente con los referentes técnicos observados en soluciones de museos inteligentes y guías por voz orientadas a accesibilidad~\cite{verde2023blemuseum,wang2024voiceguide,ivanov2025smartmuseums}.

\section{Caracterización del mercado objetivo}

El mercado objetivo de MuseIQ puede caracterizarse a partir de cinco dimensiones: institucional, tecnológica, operativa, económica y de experiencia de usuario. Desde el punto de vista institucional, se trata de organizaciones cuya misión incluye preservación, difusión, educación y valorización del patrimonio. Esto implica que sus decisiones de adopción tecnológica no responden únicamente a criterios de rentabilidad, sino también a objetivos de acceso cultural, inclusión, innovación educativa y fortalecimiento de la experiencia de visita.

Desde la dimensión tecnológica, el mercado presenta heterogeneidad. Algunos museos y espacios culturales cuentan con catálogos digitales, recorridos virtuales o plataformas web institucionales, mientras que otros mantienen procesos todavía centrados en material impreso y mediación presencial. Sin embargo, esta heterogeneidad no elimina la viabilidad de MuseIQ, ya que la solución se apoya en infraestructura relativamente ligera: beacons BLE, smartphone del visitante y una capa de servicios para contenido contextual. La existencia de plataformas como \textit{Visita Virtual} del Ministerio de Cultura demuestra que el sector ya ha iniciado procesos de digitalización cultural, aunque con foco mayor en consulta remota que en asistencia inteligente durante la visita presencial~\cite{visitavirtual2026}.

En el plano operativo, las instituciones objetivo suelen enfrentar restricciones de personal, variaciones en el flujo de visitantes y necesidades de actualización continua de contenidos. Museos con múltiples salas o circuitos requieren mecanismos que orienten al usuario y reduzcan la fricción interpretativa; museos de sitio y espacios patrimoniales necesitan recursos que ayuden a contextualizar objetos, zonas o hitos del recorrido; y centros culturales con exposiciones temporales demandan herramientas flexibles que puedan adaptarse sin rediseños complejos. En todos estos casos, MuseIQ responde con activación contextual por proximidad, entrega de narrativas en audio y consultas bajo demanda.

Desde la perspectiva económica, el mercado objetivo presenta capacidades de pago diferenciadas. Los museos públicos y sitios patrimoniales suelen depender de presupuesto estatal, fondos concursables o proyectos específicos, por lo que valoran soluciones escalables y de costo inicial contenido. Los museos privados, universitarios o con fuerte orientación de marca institucional pueden mostrar mayor disposición a invertir en herramientas que mejoren experiencia y posicionamiento. Esta diferencia sugiere una segmentación útil entre mercado primario, secundario y potencial.

\begin{table}[h!]
	\centering
	\footnotesize
	\setlength{\tabcolsep}{3pt}
	\begin{tabular}{>{\raggedright\arraybackslash}p{2.8cm} >{\raggedright\arraybackslash}p{2.8cm} >{\raggedright\arraybackslash}p{3.1cm} >{\raggedright\arraybackslash}p{3.4cm} >{\raggedright\arraybackslash}p{2.1cm}}
		\hline
		\textbf{Segmento} & \textbf{Perfil institucional} & \textbf{Necesidades predominantes} & \textbf{Aporte de MuseIQ} & \textbf{Prioridad} \\
		\hline
		Museos públicos y de sitio & Instituciones estatales con alta responsabilidad educativa y patrimonial & Mediación escalable, accesibilidad, actualización de contenidos y orientación en recorrido & Guía contextual, apoyo a recorridos autónomos y reducción de dependencia exclusiva de mediación humana & Alta \\
		Museos privados y universitarios & Instituciones con mayor flexibilidad operativa e interés en innovación & Diferenciación de experiencia, modernización y análisis de interacción del visitante & Experiencia digital autoguiada y valor agregado para públicos especializados o recurrentes & Media \\
		Centros culturales y espacios de interpretación & Entidades con exposiciones temporales o circuitos temáticos & Flexibilidad, despliegue rápido, soporte multiformato y contenidos adaptables & Configuración ligera y reutilizable para exposiciones cambiantes & Media \\
		Mercado potencial ampliado & Sitios patrimoniales abiertos y circuitos turísticos culturales & Contenido multilingüe, navegación y soporte en escenarios de mayor complejidad física & Adaptación futura a recorridos híbridos y experiencias patrimoniales extendidas & Media-baja \\
		\hline
	\end{tabular}
	\caption{Segmentación del mercado objetivo para MuseIQ. Fuente: elaboración propia con base en~\cite{museoscultura,osiptel2024smartphone,cultura2024reporte}.}
\end{table}

En términos de experiencia de usuario, el mercado objetivo comparte una necesidad transversal: ofrecer visitas más comprensibles, accesibles y personalizadas. La evidencia sobre museos inteligentes sugiere que la personalización y la interacción contextual inciden en la satisfacción y en el involucramiento del visitante~\cite{ivanov2025smartmuseums}. Asimismo, los estudios sobre guías interactivas por voz muestran valor específico para accesibilidad y autonomía en la experiencia cultural~\cite{wang2024voiceguide}. Por ello, la caracterización del mercado no debe limitarse a la institución compradora, sino incorporar al visitante como usuario final de la solución.

En función de lo anterior, el mercado primario de MuseIQ puede definirse como museos públicos, museos de sitio y espacios estatales de visita presencial con necesidad de mediación digital. El mercado secundario comprende museos privados, universitarios y centros culturales con interés en innovación de experiencia. El mercado potencial incluye espacios patrimoniales de mayor dispersión física y proyectos culturales itinerantes, cuya adopción puede ocurrir en una fase posterior una vez validado el modelo inicial.

\section{Elección del mercado objetivo}

Aunque existen varios segmentos institucionales viables, el mercado objetivo prioritario para una primera implementación de MuseIQ corresponde a museos públicos peruanos de alta o media afluencia, especialmente aquellos ubicados en Lima Metropolitana o en ciudades con ecosistemas culturales consolidados. Esta elección responde a una decisión estratégica sustentada en criterios de viabilidad técnica, impacto esperado, facilidad de validación y escalabilidad posterior.

En primer lugar, la viabilidad técnica es mayor en instituciones con circuitos definidos, espacios interiores identificables y documentación de colecciones relativamente organizada. Este contexto facilita el despliegue de beacons, la calibración de zonas de proximidad y la estructuración del repositorio documental que alimenta la capa RAG. Adicionalmente, el uso extendido de smartphones en el país favorece la interacción móvil del visitante sin exigir equipamiento adicional para cada recorrido~\cite{osiptel2024smartphone}.

En segundo lugar, la elección de museos públicos ofrece mayor impacto potencial. Las estadísticas de visitas a museos administrados por el Estado muestran recuperación sostenida y volúmenes relevantes de atención, con 1\,790\,502 visitas registradas en 2024~\cite{cultura2024reporte}. Una intervención en este segmento no solo permite validar la solución en condiciones reales de uso, sino también demostrar utilidad social y educativa en instituciones con alcance amplio y sostenido.

En tercer lugar, este segmento ofrece mejores condiciones para una validación defendible en el marco de una tesis. Un piloto en museos públicos o de sitio permite observar interacción con públicos heterogéneos, identificar necesidades de accesibilidad, analizar recorridos presenciales y producir evidencia empírica sobre aceptación, desempeño técnico y valor de uso. En contraste, iniciar únicamente en museos privados o en espacios muy pequeños podría sesgar la validación hacia contextos de mayor flexibilidad administrativa, pero con menor representatividad del ecosistema museístico nacional.

La elección prioritaria no implica descartar otros segmentos. Los museos privados y universitarios constituyen un espacio atractivo para una segunda fase debido a su potencial de adopción más ágil y a su interés en diferenciar la experiencia del visitante. No obstante, desde una perspectiva estratégica, comenzar por museos públicos de afluencia media o alta permite alinear la propuesta con necesidades sociales más amplias, generar evidencia transferible y construir una base de replicabilidad. En ese sentido, la decisión no es solamente descriptiva, sino una selección racional del segmento donde convergen necesidad real, factibilidad técnica y capacidad de demostración.

Finalmente, la baja tasa de adopción observada en aplicaciones móviles nativas de audioguías museales refuerza la conveniencia de priorizar un diseño de acceso ligero, por ejemplo mediante enlaces web o experiencias de incorporación sencilla, especialmente en el segmento público donde reducir fricción de uso es un factor crítico~\cite{sala2025museumapps}. Así, la elección del mercado objetivo también condiciona el enfoque de implementación y no solo la institución de partida.

\section{Panorama nacional}

El panorama nacional peruano muestra condiciones favorables para la adopción progresiva de una solución como MuseIQ, aunque también expone limitaciones que deben considerarse en su implementación. Desde el punto de vista institucional, el país cuenta con una red relevante de museos administrados por el Ministerio de Cultura, a la que se suman museos regionales, municipales, universitarios y privados. Esta base institucional constituye un entorno real de aplicación para soluciones de mediación cultural digital~\cite{museoscultura}.

En términos de demanda cultural, los museos estatales han mostrado una recuperación sostenida en número de visitas luego del impacto de la pandemia. La serie reciente reportada por el sector cultura confirma que el flujo presencial volvió a niveles altos y que la visita a museos mantiene relevancia como práctica cultural~\cite{cultura2020memoria,cultura2022memoria,cultura2023memoria,cultura2024reporte}. Esta recuperación fortalece la pertinencia de invertir en herramientas que acompañen la visita física y no únicamente en formatos de difusión remota.

El contexto nacional también evidencia un proceso de digitalización cultural en desarrollo. La existencia de plataformas oficiales como \textit{Visita Virtual}, así como catálogos, recorridos y contenidos digitales promovidos por entidades públicas, demuestra que el sector reconoce el valor de las tecnologías para ampliar acceso y difusión patrimonial~\cite{visitavirtual2026}. Sin embargo, estas iniciativas se concentran principalmente en visualización remota, consulta informativa o acercamiento previo y posterior a la visita, mientras que permanece abierta la oportunidad de fortalecer la mediación digital dentro del recorrido presencial.

Desde el lado del usuario, la alta penetración del smartphone es uno de los factores más relevantes para la adopción. El hecho de que 94.8\% de los hogares peruanos cuente con smartphone sugiere que una solución basada en dispositivos móviles tiene condiciones de acceso mucho más favorables que hace pocos años~\cite{osiptel2024smartphone}. No obstante, esta fortaleza debe leerse con cautela: la disponibilidad de dispositivo no garantiza por sí misma competencias digitales homogéneas, conectividad estable ni disposición a descargar aplicaciones específicas. Por ello, la implementación nacional de MuseIQ requiere interfaces sencillas, tiempos de respuesta breves y una estrategia de acceso con mínima fricción.

El panorama peruano también presenta restricciones estructurales. No todos los museos tienen el mismo nivel de infraestructura, presupuesto o madurez digital. Existen instituciones con capacidad para sostener proyectos tecnológicos y otras cuya prioridad inmediata sigue siendo la conservación básica, la señalización o la ampliación de servicios de atención al visitante. Esta heterogeneidad obliga a pensar MuseIQ como una solución modular y escalable, capaz de adaptarse a distintos grados de adopción institucional.

Aun con esas limitaciones, el contexto nacional favorece la pertinencia de la propuesta. La política cultural y las memorias de gestión del sector muestran atención sostenida al acceso, la participación y la puesta en valor del patrimonio~\cite{culturapesem2021,cultura2024reporte}. En consecuencia, una solución de guía inteligente para museos no aparece como un elemento ajeno al ecosistema, sino como una evolución de los esfuerzos de digitalización ya existentes. La oportunidad concreta consiste en cerrar la brecha entre la difusión digital del patrimonio y la asistencia contextual durante la visita física. En ese espacio se ubica MuseIQ: una herramienta orientada a mejorar orientación, interpretación y accesibilidad dentro del museo peruano contemporáneo.
